\scanpage{156}
donc est lisse à fibres connexes. D'autre part $\underline{\mathrm{SO}}(q)/\underline{\mathrm{SO}}_D(q)\simeq\underline G_{\{n\}}(q)$ est lisse à fibres géom connexes, donc aussi $\underline{\mathrm{SO}}(q)$.

\uline{Corollaire} Les $\underline{\mathrm{SO}}_D(q)$ \inserted{$=B_D(q)$} sont lisses à fibres connexes. On a $\underline{\mathrm O}_D(q)\simeq B_D(q)\times\mu_2$.

\begin{center}\rule{.12\linewidth}{.4pt}\end{center}

Supposons enfin $N=2n$, de sorte que $\mu_2=\mathrm Z(\mathrm O(q))\subset\mathrm{Sl}(E)$. On ne va pas définir $\underline{\mathrm{SO}}(q)$ comme $\mathrm O(q)\cap\mathrm{Sl}(E)$, à cause de car. $2$. Mais on fait opérer par transp de str. $\underline{\mathrm O}(q)$ sur $\underline{\mathrm{Arf}}(q)$, donc on a un hom can.
\[
\underline{\mathrm O}(q)\xrightarrow{\varepsilon_q}(\mathbb Z/2)_S
\]
\margin{explicite : soit $\eta:(\mathbb Z/2)_S\longrightarrow\mu_{2S}$ can, alors $\eta\circ\varepsilon_q=\det_q$ (restriction de $\det_E$ à $\underline{\mathrm O}(q)$) \uline{Cor} si $2$ inv. sur $S$ $\underline{\mathrm{SO}}(q)=\underline{\mathrm O}(q)\cap\mathrm{Sl}(E)$}
et on pose
\[
\underline{\mathrm{SO}}(q)=\mathrm{Ker}\,\varepsilon
\]
de sorte qu'on a la suite exacte
\[
1\mathbin{-}\underline{\mathrm{SO}}(q)\mathbin{-}\underline{\mathrm O}(q)\xrightarrow{\varepsilon_q}(\mathbb Z/2)_S\longrightarrow1
\]

\inserted{\uline{Prop} \deleted{$\mathrm{SO}(q)$} Soit $H\subset[1,n]$, \inserted{$d=\sup H$,} alors $\underline{\mathrm{SO}}(q)$ est transitif sur $\underline G_H(q)$ \deleted{\unsure{ssi|sauf si} $d\unsure{=}\sup H$} $d\ne n$, et \unsure{inv.|aussi} si $S\ne\varnothing$. Si $d=n$, et si on a une section de $\underline{\mathrm{Arf}}(S)$ sur $S$, alors $\underline{\mathrm{SO}}(q)$ est transitif sur $\underline G_H^\omega(q)$.}%
% margin note written sideways

Considérons comme dessus un $D$, d'où
% archive sheet number: 155
